% =========================================================
% Chapter: Structure of the Real Line
% =========================================================
\chapter{Structure of the Real Line}
\label{chap:structure-of-real-line}

\breadcrumb{structure-of-real-line}{Number Lines and Intervals}{Structure of the Real Line}{Formalizing the Number Systems}


\begin{tcolorbox}[
  colback=gray!6,
  colframe=gray!40,
  arc=2pt,
  left=8pt, right=8pt, top=6pt, bottom=6pt,
  title={\small\textbf{Structural Roadmap}},
  fonttitle=\small\bfseries
]
\noindent\textbf{How we got here.}
The embedding chapter placed $\mathbb{N}$, $\mathbb{W}$, $\mathbb{Z}$, and $\mathbb{Q}$ inside $\mathbb{R}$, and the number-line chapter fixed the interval notation used to describe subsets of the line.

\medskip
\noindent\textbf{What this chapter develops.}
This chapter develops only the metric topology of $\mathbb{R}$: the distance $d(x,y)=|x-y|$, balls as open intervals, neighborhoods, open and closed sets, interior and boundary, limit points, covers and finite subcovers, compactness by the open-cover definition, and the Heine--Borel theorem as capstone. It is deliberately not a general topology chapter.

\medskip
\noindent\textbf{Where this leads.}
Heine--Borel certifies that the closed bounded sets are exactly the compact ones, which is the entry point to Analysis~I; the open-cover definition of compactness is the form that later carries to the torus, where ``bounded'' has no intrinsic meaning.
\end{tcolorbox}

% =========================================================
% Structure of the Real Line — Notes Index
% =========================================================
\section{Ordered Field Structure}

\input{volume-ii/structure-of-real-line/notes/ordered-field-structure/notes-ordered-field-structure}


\input{volume-ii/structure-of-real-line/notes/density-and-no-adjacent-points/notes-density-and-no-adjacent-points}


\input{volume-ii/structure-of-real-line/notes/intervals-as-ordered-subsets/notes-intervals-as-ordered-subsets}


\section{Completeness as Real Line Feature}

\input{volume-ii/structure-of-real-line/notes/completeness-as-real-line-feature/notes-completeness-as-real-line-feature}


\input{volume-ii/structure-of-real-line/notes/from-number-systems-to-analysis/notes-from-number-systems-to-analysis}




\section*{Proofs}
\LRAProofsInput{volume-ii/structure-of-real-line/proofs/index}

\section*{Capstone}
\LRAExercisesInput{volume-ii/structure-of-real-line/proofs/exercises/index}
